% Content of the blueprint for Chapter 1 of
% Peter Gacs, "Descriptional Complexity and Randomness" (arXiv:2105.04704).
%
% Statements and proofs are his, reproduced with the author's permission.
% Each node carries the name of the corresponding Lean declaration.  Where our
% formalized statement departs from the printed one, a formalization note says so.
%
% This file is assembled to keep his numbering: within each section, theorems run
% on one counter and everything else on another, exactly as in the manuscript.

\chapter{Complexity}

This blueprint follows Chapter 1 of Peter G\'acs, \emph{Descriptional Complexity and
Randomness}~\cite{Gacs2021}, and records how each of its results is formalized in Lean 4
on top of Mathlib. The statements and their proofs are his, reproduced with his
permission; the annotations, the dependency edges and the formalization notes are ours.

Every numbered result carries the number it has in his manuscript, so the two can be read
side by side. A green node is formalized. Where the Lean statement is not a literal
transcription of the printed one, a \emph{formalization note} follows it and says what
changed and why. Two topics are stated here without being formalized: the clauses of
Proposition 1.5.5 about functions of a real variable, and Theorem 1.6.4 on exact
domination.

\section{Introduction}

His first two sections are informal. Section 1.1 motivates description complexity and
surveys the results proved later, and Section 1.2 fixes notation. Neither contains a
result that the formalization has to reproduce, so they appear here only to keep the
section numbering aligned with the manuscript. The notation we do rely on is recalled
below.

\section{Notation }

We write $\dS$ for the set of finite binary strings and $\len{x}$ for the length of a
string $x$. Strings are paired by a computable encoding, written $\ang{x,y}$, and
$J(n) = n + 2\log n$ is the cost of making a string of length $n$ self-delimiting. The
relation $f \lea g$ means that $f \le g + c$ for a constant $c$ that does not depend on
the arguments, and $f \eqa g$ means that both $f \lea g$ and $g \lea f$ hold. Their
multiplicative counterparts $\lem$ and $\eqm$ mean the same up to a constant factor.

In Lean, a string is a \texttt{List Bool}, written \texttt{BitString}, and a decompressor
is a partial recursive map \texttt{Map}. Complexities are valued in $\dN \cup \setof{\infty}$,
and every quantity that his text writes as a negative logarithm is carried multiplicatively
in $[0,\infty]$ instead.



\section{Kolmogorov complexity }
The complexity of a string is the length of its shortest description under a fixed
optimal interpreter. The invariance theorem is what makes the definition well posed, and
everything later in the chapter rests on it.


% Gacs source line 796
 \begin{definition}
  \label{def:interpreter}
  \lean{Kolmogorov.Map}
  \leanok
A partial recursive 
function \( A \) from \( \dS_{2} \times \dS \) to \( \dS \) will be called a (binary)
\df{interpreter}, or \df{decoder}. 
We will often refer to its first argument as the \df{program}, or \df{code}. 
 \end{definition}

% Gacs source line 807
 \begin{definition}
  \label{def:complexity}
  \lean{Kolmogorov.condK, Kolmogorov.plainK}
  \leanok
  \uses{def:interpreter}
For any binary interpreter \( A \) and strings \( x,y \in \dS \), the number
 \[
   C_A(x \mvert y) = \min \setof{ \len{p} : A(p,y) = x }
 \] 
is called the \df{conditional Kolmogorov-complexity of \( x \) given
\( y \), with respect to the interpreter \( A \) }. 
(If the set after the ``min'' is empty, then the minimum is \( \infty \)). 
We write  \( C_A(x)=C_A(x \mvert \Lg) \).  
 \end{definition}

% Gacs source line 830
 \begin{definition}
  \label{def:optimal}
  \lean{Kolmogorov.isOptimalConditional}
  \leanok
  \uses{def:interpreter, def:complexity}
A binary p.r.~interpreter \( U \) is called \df{optimal}
if for any binary p.r.~interpreter \( A \) there is a constant \( c < \infty \) such that 
for all \( x,y \) we have
 \begin{equation}  
C_U(x \mvert y) \le C_A(x \mvert y) + c. \label{eq:opt}
 \end{equation}
 \end{definition}

% Gacs source line 839
\begin{theorem}[Invariance Theorem]
  \lean{Kolmogorov.existsIsOptimalConditional}
  \leanok
  \uses{def:optimal} \label{thm:invar}
 There is an optimal p.r.~binary interpreter.
 \end{theorem}

% Gacs source line 842
\begin{proof}
  \leanok
  \uses{def:optimal} 
The idea is to use interpreters that come from universal partial recursive
functions.
However,  not all such functions can be used for universal interpreters.
Let us introduce an appropriate pairing function.
For \( x \in \dB^{n} \), let 
 \[
   x^{o} = x(1) 0 x(2) 0 \ldots x(n-1) 0 x(n) 1 
 \]
  where \( x^o = x1 \) for \( \len{x} =1 \). Any binary sequence \( x \) can be
uniquely represented in the form \( p = a^o b \). 

We know there is a p.r.~function
\( V : \dS_{2} \times \dS_{2} \times \dS \to \dS \) that
is \emph{universal}: for any
p.r.~binary interpreter \( A \), there is a string \( a \) such that for all
\( p,x \), we have \( A(p,x) = V(a,p,x) \). 
Let us define the function \( U(p,x) \) as follows. 
We represent the string \( p \) in the form 
\( p = u^o v \) and define \(  U(p,x) = V(u,v,x)  \). 
Then \( U \) is a p.r.~interpreter.
Let us verify that it is optimal. 
Let \( A \) be a p.r.~binary
interpreter, \( a \) a binary string such that \( A(p,x) = U(a,p,x) \) for
all \( p,x \). 
Let \( x,y \) be two strings. 
If \( C_A(x \mvert y)=\infty \), then~\eqref{eq:opt} holds trivially. 
Otherwise, let \( p \) be a binary string of
length \( C_A(x \mvert y) \) with \( A(p,y)=x \). 
Then we have
 \[
  U(a^o p, y) = V(a,p,y) = A(p,y) = x, 
 \]
  and
 \[
  C_U(x \mvert y) \le 2\len{a} + C_A(x \mvert y).
 \]
  \end{proof}

% Gacs source line 901
 \begin{definition}
  \label{def:fixed-U}
  \uses{thm:invar, def:complexity}
We fix an optimal binary p.r.~interpreter \( U \) and write 
\( C(x \mvert y) \) for \( C_U(x \mvert y) \). 
 \end{definition}

% Gacs source line 919
 \begin{definition}
  \label{def:tuples}
  \uses{def:fixed-U}
We define \( U_{p}(x_{1},\ldots,x_{k}) = U_{p}(\ang{ x_{1}\,\ldots, x_{k}}) \). 
Similarly,
we will refer to an arbitrary computably enumerable
set as the range \( W_{p} \) of some \( U_{p} \).  
We often do not need the second argument of the function \( U(p,x) \).  
We therefore define 
 \[ U(p) = U(p,\Lg).
 \]
 \end{definition}

% Gacs source line 954
 \begin{definition}
  \label{def:additive-order}
The relation 
\( f \lea g \) means inequality to within an additive constant, that
there is a constant \( c \) such that for all \( x \), \( f(x) \le g(x) \). 
We can write this also as \( f \le g + O(1) \). 
We also say that \( g \) \df{additively dominates} \( f \). 
The relation \( f \eqa g \) means \( f \lea g \)
and \( f \gea g \). 
The relation \( f \lem g \) among nonnegative functions
means inequality to within a multiplicative constant. 
It is
equivalent to \( \log f \lea \log g \) and \( f = O(g) \). 
We also say that
\( g \) \df{multiplicatively dominates} \( f \). 
The relation \( f\eqm g \) means \( f \lem g \) and \( f \gem g \).
 \end{definition}

% Gacs source line 973
\begin{theorem}
  \lean{Kolmogorov.plainKNatLeLength, Kolmogorov.plainKLeLength, Kolmogorov.cardCompressibleWordsLt}
  \leanok
  \uses{def:fixed-U, def:additive-order} \label{thm:upbd}
The following simple upper bound holds.
 \begin{alphenum}  
  \item \label{thebd} For any natural number \( m \), we have 
 \begin{equation} C(m) \lea \log m. \label{eq:upbd}
 \end{equation}
 \item \label{howsharp}
  For any positive real number \( v \) and string \( y \), every finite set 
	\( E \) of size \( m \) has at least \( m(1-2^{-v+1}) \) elements \( x \) with 
	\( C(x \mvert y) \ge \log m - v  \).
 \end{alphenum}
\end{theorem}

% Gacs source line 986
\begin{corollary}
  \label{cor:limC}
  \lean{Kolmogorov.plainKNat_eventually_gt, Kolmogorov.plainKNat_toNat_tendsto_atTop}
  \leanok
  \uses{thm:upbd}  \( \lim_{n\to\infty} C(n) = \infty \).
\end{corollary}

% Gacs source line 993
\begin{proof}[Proof of Theorem \protect\ref{thm:upbd}]
  \proves{thm:upbd}
  \leanok
  \uses{def:fixed-U}
 First we prove~\eqref{thebd}.  Let the interpreter \( A \) be defined such
that if \( \bg(m) = p \) then \( A(p,y) = m \). Since
 \( |\bg(m)| = \cei{ \log m} \), we have \( C(m) \lea C_A(m) < \log m + 1 \).

Part~\eqref{howsharp} says that the trivial estimate~\eqref{eq:upbd} is
quite sharp for {\em most} numbers under \( m \). The reason for this is
that since there are only few short programs, there are only few
objects of low complexity. For any string \( y \), and any positive real
natural number \( u \), we have 
 \begin{equation} \card{\setof{ x : C(x \mvert y) \le u }} < 2^{u+1}. 
\label{eq:howsharp} \end{equation} 
  To see this, let \( n=\flo{ \log u } \). The number \( 2^{n+1}-1 \) of 
different binary strings of length \( \le n \) is an upper 
bound on the number of different shortest programs of length \( \le u \). 
Now~\eqref{howsharp} follows immediately from~\eqref{eq:howsharp}.
 \end{proof}

%%%% ===================== SECTION 1.4 Simple properties of information =====================


\section{Simple properties of information }
The results of this section are all consequences of one move: given an interpreter, build
another one that does a little extra work before or after it, then appeal to optimality.
In Lean the same move appears as a pair of transport lemmas.


% Gacs source line 1071
\begin{theorem}
  \lean{Kolmogorov.condK_le_condK_ctx_map}
  \leanok
  \uses{def:fixed-U, def:additive-order} \label{thm:rectransf}
 For any partial recursive function \( U_{q} \), over 
strings, we have
 \begin{align*} C(U_{q}(x) \mvert y)	&\lea  C(x \mvert y) + J(\len{q}),
  \\   C(y \mvert U_{q}(x))	&\gea  C(y \mvert x) - J(\len{q}).
 \end{align*}
\end{theorem}

% Gacs source line 1079
\begin{proof}
  \leanok
  \uses{def:optimal}
 To define an interpreter \( A \) with \( C_A(U_{q}(x) \mvert y) \lea C(x \mvert y )
+ J(\len{q}) \), we define \( A(\ig(q)p,y) = U_{q}(U(p,y)) \). 
To define an
interpreter \( B \) with \( C_B(y \mvert x) \lea C(y \mvert U_{q}(x)) + J(q) \), we
define \( B(\ig(q)p,x) = U(p,U_{q}(x)) \).
 \end{proof}

% Gacs source line 1089
 \begin{definition}
  \label{def:pair-complexity}
  \uses{def:tuples, def:fixed-U}
The definition of conditional complexity is extended to pairs, etc. by 
 \[
 C(x_{1},\ldots,x_{m} \mvert y_{1},\ldots, y_{n}) 
	= C(\ang{ x_{1},\ldots,x_{m}}  \mvert \ang{ y_{1},\ldots,y_{n}} ).
 \]
 \end{definition}

% Gacs source line 1099
\begin{corollary}
  \lean{Kolmogorov.condK_le_condK_pair_left, Kolmogorov.condK_le_condK_pair_right, Kolmogorov.plainK_pair_self_eq}
  \leanok
  \uses{thm:rectransf, def:pair-complexity} \label{c:upbdrec}
 For any one-to-one p.r. function \( U_{p} \), we have 
 \[ |C(x) - C(U_{p}(x))| \lea J(\len{p}).
 \]
  Further,
 \begin{align} 	C(x \mvert y,z)	&\lea  C(x \mvert U_{p}(y),z)+J(\len{p}), \nn
  \\	C(U_{p}(x)) 	&\lea  C(x) + J(\len{p}),\nn	
  \\	C(x \mvert z)	&\lea  C(x,y \mvert z), \label{eq:x.xy} 		
  \\	C(x \mvert y,z)	&\lea  C(x \mvert y),   \nn		
  \\	C(x,x)		&\eqa  C(x),       \label{eq:xx.x}
  \\	C(x,y \mvert z)	&\eqa  C(y,x \mvert z),	\nn	
  \\	C(x \mvert y,z)	&\eqa  C(x \mvert z,y),	\nn	
  \\	C(x,y \mvert x,z)	&\eqa  C(y \mvert x,z),	\nn	
  \\	C(x \mvert x,z)	&\eqa  C(x \mvert x) \eqa 0.  \nn		
 \end{align}
 \end{corollary}

% Gacs source line 1124
 \begin{theorem}
  \lean{Kolmogorov.plainK_pairCode_le}
  \leanok
  \uses{def:pair-complexity} \label{thm:addit}
 \[  C(x,y) \lea J(C(x)) + C(y \mvert x).
 \]
 \end{theorem}

% Gacs source line 1129
\begin{proof}
  \leanok
  \uses{def:optimal, def:pair-complexity}
 We define an interpreter \( A \) as follows. 
We decompose any binary
string \( p \) into the form \( p = \ig(w)q \), and let \( A(p,z) = U(q,U(w)) \).
Then \( C_A(x,y) \) is bounded, to within an additive constant, by the 
right-hand side of the above inequality.
 \end{proof}

% Gacs source line 1137
\begin{corollary}
  \lean{Kolmogorov.plainK_pairCode_le_plainK, Kolmogorov.plainK_map_pairCode_le}
  \leanok
  \uses{thm:addit, thm:rectransf} \label{cy:funcofpair}
 For any p.r.~function \( U_{p} \) over pairs of strings, we have 
 \begin{align*} 
C(U_{p}(x,y)) &\lea  J(C(x)) + C(y \mvert x) + J(\len{p})
  \\		  &\lea  J(C(x)) + C(y) + J(\len{p}),
 \end{align*}
  and in particular,
 \[
 	C(x,y) \lea J(C(x)) + C(y).
 \]
 \end{corollary}

% Gacs source line 1205
\begin{theorem}
  \lean{Kolmogorov.condK_le_log_card, Kolmogorov.condK_le_log_card_finset}
  \leanok
  \uses{def:pair-complexity} \label{thm:condupbd}
 Let \( E=W_{p} \) be an enumerable set of pairs of strings defined
enumerated with the help of the program \( p \) for example as 
 \begin{equation} 
 W_{p}= \setof{ U(p,x): x\in \dN}. \label{eq:Wp}
 \end{equation}
  We define the \df{section} 
 \begin{equation} 
 E^a = \setof{ x :  \tupcod{a,x}  \in E }. \label{eq:sectiondef}
 \end{equation}
  Then for all \( a \) and \( x \in E^a \), we have 
 \[  
  C(x \mvert a) \lea \log\card{E^a} + J(\len{p}).
 \]
 \end{theorem}

% Gacs source line 1220
\begin{proof}
  \leanok
  \uses{def:optimal}
 We define an interpreter \( A(q,b) \) as follows. 
We decompose \( q \) into \( q = \ig(p) \bg(t) \). 
From a standard enumeration \( (a(k),x(k)) \)
\( (k=1,2,\ldots) \) of the set \( W_{p} \), we produce an enumeration \( x(k,b) \)
\( (k=1,2,\ldots) \), without repetition, of the set \( W_{p}^b \) for each
\( b \), and define \( A(q,b)=x(t,b) \). 
For this interpreter \( A \), we get
\( k_A(x \mvert a) \lea J(\len{p}) + \log \card{E^a} \). 
 \end{proof}

% Gacs source line 1239
\begin{theorem}
  \lean{Kolmogorov.plainK_pair_selfComplexity_eq}
  \leanok
  \uses{def:pair-complexity, thm:upbd} \label{thm:xKx}
We have \( C(x,C(x)) \eqa C(x) \).
 \end{theorem}

% Gacs source line 1242
\begin{proof}
  \leanok
  \uses{def:optimal}
 The inequality \( \gea \) follows from~\eqref{eq:x.xy}. 
To prove \( \lea \), let \( A(p) =  \ang{ U(p,\len{p})}  \). 
Then \( C_A(x,C(x)) \le C(x) \).
 \end{proof}

% Gacs source line 1250
\begin{theorem}
  \lean{Kolmogorov.condK_pair_sub_selfComplexity_le}
  \leanok
  \uses{def:pair-complexity} \label{thm:strange}
 \[	
C(y\mvert x,\ i-C(y\mvert x,i)) \lea C(y\mvert x,i).
 \]
 \end{theorem}

%%%% ===================== SECTION 1.5 Algorithmic properties of complexity =====================


\section{Algorithmic properties of complexity}
Complexity is not computable, but it is upper semicomputable, and that one-sided
computability is exactly what makes the theory work. Levin's theorem turns the
observation into a characterization.


% Gacs source line 1269
\begin{definition}[Computability]
  \label{def:computable-function}
  \lean{Kolmogorov.IsComputableENNReal}
  \leanok
Let \( f:\dS\to\dR \) be a function.
It is \df{computable} if there is a recursive function \( g(x,n) \) with rational
values, and  \( |f(x)-g(x,n)| < 1/n \). 
\end{definition}

% Gacs source line 1275
 \begin{definition}[Semicomputability]
  \lean{Kolmogorov.IsLSC₁, Kolmogorov.UpperSemicomputable}
  \leanok
  \uses{def:computable-function}\label{def:semicomputability}
A function \( f : \dS \to \rint{-\infty}{\infty} \) is \df{(lower) semicomputable}
if the set 
 \[
	\setof{ \tup{x,r} : x \in \dS,\; r \in \dQ,\; r<f(x) }
 \] 
 is recursively enumerable. 
A function \( f \) is called \df{upper semicomputable} if \( -f \) is
lower semicomputable.
 \end{definition}

% Gacs source line 1305
\begin{definition}[Computability for real functions]
  \label{def:computable-real-function}
  \notready
Let \( f:\dR\to\dR \) be a function.
It is \df{computable} if there is a recursively enumerable set 
\( \cF\sbsq \bg^{2} \)
such that denoting \( \cF_{J}=\setof{I: (I,J)\in \cF} \) we have
 \begin{align*}
   f^{-1}(J) = \bigcup_{I\in\cF_{J}}I.
 \end{align*}
\end{definition}

\begin{formalization}
This clause and the next concern functions of a real variable, in the sense of enumerable sets of rational intervals. The library has no computability theory for real-input functions, and nothing else in Chapter 1 depends on one, so these two definitions are stated here but not formalized.
\end{formalization}

% Gacs source line 1318
\begin{definition}
  \label{def:lsc-real-function}
  \notready
  We say that \( f:\dR\to\ol\dR \) is
\df{lower semicomputable} if there is a recursively enumerable set 
\( \cG\sbsq \cQ\times\bg \)
such that denoting \( \cG_{q}=\setof{I: (I,q)\in \cG} \) we have
 \begin{align*}
   f^{-1}(\opint{q}{\infty}) = \bigcup_{I\in\cF_{J}}I.
 \end{align*}
\end{definition}

% Gacs source line 1332
\begin{proposition}
  \lean{Kolmogorov.IsLSC₁.add, Kolmogorov.IsLSC₁.max, Kolmogorov.IsLSC₁.comp_computable, Kolmogorov.IsComputableENNReal.comp_computable}
  \uses{def:computable-function, def:semicomputability}\label{propo:l-sc-properties}\ 
  \begin{alphenum}
  \item 
  The function \( \cei{\cdot}:\dR\to\dR \) is lower semicomputable.
  \item If \( f,g:\dR\to\dR \) are computable then their composition
    \( f(g(\cdot)) \) is, too.
  \item If \( f:\dR\to\dR \) is lower semicomputable and \( g:\dR\to\dR \) 
(or \( \dS\to\dR \)) is computable then \( f(g(\cdot)) \) is lower semicomputable.
  \item If \( f:\dR\to\dR \) is lower semicomputable and monotonic,
and \( g:\dR\to\dR \) (or \( \dS\to\dR \)) is lower 
semicomputable then \( f(g(\cdot)) \) is lower semicomputable.
  \end{alphenum}
\end{proposition}

\begin{formalization}
Clauses (b) and (c) are formalized for functions on strings, together with closure under sum and maximum. Clauses (a) and (d) speak about functions of a real variable and are not formalized, which is why this node is not marked green.
\end{formalization}

% Gacs source line 1346
 \begin{definition}
  \label{def:universal-lsc}
  \uses{def:semicomputability}
We introduce a \df{universal lower semicomputable function}. 
For every binary string \( p \), and \( x \in \dS \), let
 \[ 
  S_{p}(x)= \sup\setof{ y/z : y,z\in Z, z \not = 0, 
	\ang{\ang{ x } ,y,z }  \in W_{p} }
 \]
where \( W_{p} \) was defined in~\eqref{eq:Wp}.  
 \end{definition}

% Gacs source line 1360
 \begin{definition}
  \label{def:lsc-tuples}
  \uses{def:universal-lsc}
Let \( S_{p}(x_{1},\ldots,x_{k})= S_{p}(\ang{ x_{1},\ldots,x_{k} } ) \).  
For any lower semicomputable function \( f \), we call any binary string \( p \) with
\( S_{p}=f \) a \df{G\"odel number} of \( f \). 
There is no universal computable function, but for any computable 
function \( g \), we call any number \(  \tupcod{  \tupcod{ p } , \tupcod{ q } }  \) 
a G\"odel number of \( g \) if \( g=S_{p}=-S_{q} \).
 \end{definition}

% Gacs source line 1371
\begin{theorem}
  \lean{Kolmogorov.condK_upperSemicomputable}
  \leanok
  \uses{def:semicomputability, thm:upbd}	      \label{thm:Ksemicp}
 The function \( C(x \mvert y) \) is upper semicomputable.
 \end{theorem}

% Gacs source line 1374
\begin{proof}
  \leanok
  \uses{thm:upbd} 
 By Theorem~\ref{thm:upbd} there is a constant \( c \) such that for
any strings \( x,y \), we have \( C(x \mvert y) < \log \tupcod{  x  }  + c \).  
Let some Turing machine compute our optimal binary p.r.~interpreter. 
Let \( U^{t}(p,y) \) be defined as \( U(p,y) \) if this machine, when started on
input \( (p,y) \), gives an output in \( t \) steps, undefined otherwise. 
Let \( K^{t}(x \mvert y) \) be the smaller of \( \log  \tupcod{  x  }  + c \) and 
 \[  
   \min\setof{ \len{p} \le t : U^{t}(p,y)=x }.
 \]
Then the function \( K^{t}(x \mvert y) \) is computable, monotonic in \( t \) and 
\( \lim_{t\to\infty} K^{t}(x \mvert y) = C(x \mvert y) \). 
 \end{proof}

% Gacs source line 1391
\begin{theorem}
  \lean{Kolmogorov.noComputableUnboundedLowerBound, Kolmogorov.notComputablePlainKNat}
  \leanok
  \uses{def:fixed-U} \label{thm:nolb}
 Let \( U_{p}(n) \) be a partial recursive function whose values are numbers
smaller than \( C(n) \) whenever defined. Then 
 \[  
 U_{p}(n) \lea J(\len{p}).
 \]
 \end{theorem}

% Gacs source line 1398
\begin{proof}
  \leanok
  \uses{def:fixed-U}
 The proof of this theorem resembles ``Berry's paradox '', which says:
``The least number that cannot be defined in less than 100 words'' is a 
definition of that number in 12 words. 
The paradox can be used to 
prove, whenever we agree on a formal notion of ``definition'', that the 
sentence in quotes is not a definition. 

Let \( x(p,k) \) for \( k=1,2,\ldots \) be an enumeration of the domain of 
\( U_{p} \). For any natural number \( m \) in the range of \( U_{p} \), let \( f(p,m) \) be 
equal to the first \( x(p,k) \) with \( U_{p}(x(p,k)) = m \). Then by definition, 
\( m < C(f(p,m)) \). On the other hand, \( f(p,m) \) is a p.r.~function, hence 
applying Corollary~\ref{cy:funcofpair} and~\eqref{eq:binupbd} we get 
 \begin{align*}
     m	&< C(f(p,m)) \lea C(p) + J(C(m))
   \\	&\lea \len{p} + 1.5\log m. 		
 \end{align*}
 \end{proof}

% Gacs source line 1427
 \begin{definition}
  \label{def:simple-set}
A computably enumerable set is \df{simple} if its complement is infinite but
does not contain any infinite computably enumerable subsets. 
 \end{definition}

% Gacs source line 1464
\begin{corollary}
  \label{cor:chaitin}
  \lean{Kolmogorov.FormalSystem.chaitinBound, Kolmogorov.FormalSystem.chaitinIncompleteness}
  \leanok
  \uses{thm:nolb}
There is a constant \( c<\infty \) such that for any \( p \), if any sentences
with the meaning ``\( m < C(n) \)'' for \( m > J(\len{p}) + c \)  are 
provable in theory \( T_{p} \) then some of these sentences are false. 
 \end{corollary}

\begin{formalization}
The Lean version states the bound and the incompleteness consequence for an abstract sound formal system, in the style of Kritchman and Raz, rather than for a concrete axiomatization of arithmetic.
\end{formalization}

% Gacs source line 1469
\begin{proof}
  \proves{cor:chaitin}
  \leanok
  \uses{thm:nolb} For some \( p \), suppose that all sentences ``\( m<C(n) \)''
provable in \( T_{p} \) are true. 
For a sentence \( \Phi \), let \( T_{p} \vdash \Phi \) 
denote that \( \Phi \) is a theorem of the theory \( T_{p} \). 
Let us define the function 
 \[
 A(p,n) = \max\setof{ m : T_{p} \vdash \text{``}m < C(n)\text{''} }.
 \]
 This function is semicomputable since the set
\( \setof{(p,\tupcod{\Phi}): T_{p}\vdash\Phi} \) is recursively
enumerable.  
Therefore there is a binary string \( a \) such that
\( A(p,n)=S(a,\tupcod{ p,n}) \). 
It is easy to see that we have a constant string \( b \) such that 
 \[
   S(a, \tupcod{ p,n } ) = S(q, n)
 \]
  where \( q = \ol b \ol a p \). 
(Formally, this statement follows for example from the so-called \( S_{m}^{n} \)-theorem.)  
The function \( S_{q}(n) \) is a
lower bound on the complexity function \( C(n) \). 
By an immediate
generalization of Theorem~\ref{thm:nolb} for semicomputable functions),
we have \( S(q,n) \lea J(\len{q}) \lea  J(\len{p}) \).
 \end{proof}

% Gacs source line 1503
\begin{theorem}[Levin]
  \lean{Kolmogorov.levin_characterization}
  \leanok
  \uses{def:semicomputability, thm:upbd}  \label{thm:Kmajor}
Let \( F(x,y) \) be a function of strings semicomputable from above. 
The relation \( C(x \mvert y) \lea F(x,y) \) holds for all \( x,y \) if and only if 
we have 
 \begin{equation} \label{eq:Kmajor}
  \log |\setof{ x : F(x,y ) < m } | \lea m  
 \end{equation}
  for all strings \( y \) and natural numbers \( m \). 
 \end{theorem}

% Gacs source line 1512
\begin{proof}
  \leanok
  \uses{thm:upbd}
 Suppose that \( C(x \mvert y) \lea F(x,y) \) holds. 
Then~\eqref{eq:Kmajor} follows from~\ref{howsharp} of Theorem~\ref{thm:upbd}.

Now suppose that~\eqref{eq:Kmajor} holds for all \( y,m \). 
Then let \( E \) be the computably enumerable set of triples \( (x,y,m) \) with \( F(x,y) < m \). 
It follows from~\eqref{eq:Kmajor} and Theorem~\ref{thm:condupbd} that for all
\( (x,y,m)\in E \) we have \( C(x \mvert y,m) \lea m \).  
By Theorem~\ref{thm:strange} then \( C(x \mvert y) \lea m \). 
 \end{proof}

%%%% ===================== SECTION 1.6 The coding theorem =====================


\section{The coding theorem}
Restricting to self-delimiting interpreters gives prefix complexity $K$, and the coding
theorem identifies it with the negative logarithm of the universal semimeasure. This is
the technical heart of the chapter.


% Gacs source line 1547
 \begin{definition}
  \label{def:prefix-free}
A set of strings is called \df{prefix-free} if for any pair \( x,y \) of
elements in it, \( x \) is not a prefix of \( y \).

A one-to-one function into a set of strings is a
\df{prefix code}, or \df{instantaneous code} if its domain of definition is
prefix free. 

An interpreter
\( f(p,x) \) is \df{self-delimiting (s.d.)} if for each \( x \), the set
\( D_{x} \) of strings \( p \) for which \( f(p,x) \) is defined is a prefix-free set.
 \end{definition}

% Gacs source line 1560
 \begin{example}
  \label{ex:selfdelim-code}
  \uses{def:prefix-free}
The mapping \( p \to p^o \) is a prefix code.
 \end{example}

% Gacs source line 1567
\begin{definition}
  \label{def:sd-machine}
  \lean{Kolmogorov.IsPrefixMachine}
  \leanok
  \uses{def:prefix-free}
A Turing machine \( \cT \) is \df{self-delimiting} if it
has no input tape, but can ask any time for a new input symbol. 
After some computation and a few such requests (if ever) the
machine decides to write the output and stop. 
\end{definition}

% Gacs source line 1589
\begin{definition}
  \label{def:sd-optimal}
  \lean{Kolmogorov.IsOptimalPrefixConditional}
  \leanok
  \uses{def:sd-machine}
A self-delimiting partial recursive 
interpreter \( T \) is called \df{optimal} if for any other s.d.~p.r.~interpreter \( F \), 
we have 
 \begin{equation} 	C_{T} \lea C_{F}. \label{eq:sdopt}  
 \end{equation}
\end{definition}

% Gacs source line 1597
\begin{theorem}
  \lean{Kolmogorov.exists_isOptimalPrefixConditional}
  \leanok
  \uses{def:sd-optimal} \label{thm:sdopt} There is an optimal s.d.~interpreter.
\end{theorem}

% Gacs source line 1599
\begin{proof}
  \leanok
  \uses{def:sd-optimal, thm:invar}								
The proof of this theorem is similar to the proof of Theorem~\ref{thm:invar}.
We take the universal partial recursive function \( V(a,p,x) \). 
We transform each function 
\( V_a(p,x)=V(a,p,x) \) into a self-delimiting function
\( W_a(p,x)=W(a,p,x) \), but so as not to change the functions \( V_a \)
which are already self-delimiting. 
Then we form \( T \) from \( W \) just as we formed \( U \) from \( V \). 
It is easy to check that the function \( T \)
thus defined has the desired properties. 
Therefore we are done if we construct \( W \).

Let some Turing machine \( \cM \) compute \( y=V_a(p,x) \) in \( t \) steps.
We define \( W_a(p,x)=y \) if \( \len{p} \le t \) and if \( \cM \) does not
compute in \( t \) or fewer steps any output \( V_a(q,x) \) from any
extension or prefix \( q \) of length \( \le t \) of \( p \). 
If \( V_a \) is
self-delimiting then \( W_a=V_a \). But \( W_a \) is always self-delimiting.
Indeed, suppose that \( p_{0} \prefix p_{1} \) and that \( \cM \) computes
\( V_a(p_{i},x) \) in \( t_{i} \) steps. 
The value \( W_a(p_{i},x) \) can be defined only if \( \len{p_{i}} \le t_{i} \). 
The definition of \( W_a \) guarantees that if
\( t_{0} \le t_{1} \) then \( W_a(p_{1},x) \) is undefined, otherwise \( W_a(p_{0},x) \)
is undefined. 
 \end{proof}

% Gacs source line 1625
 \begin{definition}
  \label{def:K}
  \lean{Kolmogorov.KP, Kolmogorov.KPPlain}
  \leanok
  \uses{thm:sdopt}
We fix an optimal self-delimiting p.r.~interpreter \( T(p,x) \) and write
 \[ 	
K(x \mvert y) = C_{T}(x \mvert y).
 \]
We call \( K(y \mvert x) \) the (self-delimiting) \df{conditional complexity} 
of \( x \) with respect to \( y \). 
 \end{definition}

% Gacs source line 1653
\begin{theorem}
  \lean{Kolmogorov.plainK_le_KPPlain_le_plainK_add_log, Kolmogorov.KP_le_condK_add_log, Kolmogorov.KPPlain_le_plainK_add_log}
  \leanok
  \uses{def:K, def:complexity} \label{thm:HK} 
 \begin{equation} K \lea H \lea J(K). 
 \end{equation}
 \end{theorem}

% Gacs source line 1657
\begin{proof}
  \leanok
  \uses{def:K} Obviously, \( K \lea H \). We define a self-delimiting
p.r.~function \( F(p,x) \). 
The machine computing \( F \) tries to decompose 
\( p \) into the form \( p=u^o v \) such that \( u \) is the number \( \len{v} \) in binary
notation. 
If it succeeds, it outputs \( U(v,x) \). 
We have 
 \[ 
K(y \mvert x) \lea C_{F}(y \mvert x) \lea C(y \mvert x) + 2\log C(y \mvert x).
 \]
 \end{proof}

% Gacs source line 1707
  \begin{definition}
  \label{def:semimeasure}
  \lean{Kolmogorov.IsSemimeasure, Kolmogorov.IsLowerSemicomputableSemimeasure}
  \leanok
  \uses{def:semicomputability}
A function \( w : S \to \clint{0}{1} \) is called a \df{semimeasure}
over the space \( S \) if 
 \begin{equation} \label{eq:dscrsemi} 
  \sum_{x} w(x) \le 1. 
 \end{equation}
 It is called a \df{measure} (a probability distribution) if
equality holds here. 

A semimeasure is called \df{constructive} if it
is lower semicomputable. 
A constructive semimeasure which
multiplicatively dominates every other constructive semimeasure is
called \df{universal}.
  \end{definition}

% Gacs source line 1723
 \begin{remark}
  \label{rem:semimeasure-later}
Just as in an earlier remark, we warn that semimeasures will later be
defined over the space \( \dN^{\dN} \). 
They will be characterized by a
nonnegative function \( w \) over \( \dS=\dN^{*} \) but with a condition different
from~\eqref{eq:dscrsemi}.
 \end{remark}

% Gacs source line 1733
 \begin{proposition}
  \label{prop:constructive-measure}
  \lean{Kolmogorov.isComputable_of_isLSC₁_of_tsum_eq_one, Kolmogorov.IsLowerSemicomputableSemimeasure.isComputable_of_tsum_eq_one}
  \leanok
  \uses{def:semimeasure}
If a constructive semimeasure \( w \) is also a measure then it is
computable. 
 \end{proposition}

% Gacs source line 1737
 \begin{proof}
  \proves{prop:constructive-measure}
  \leanok
  \uses{def:semimeasure}
If we compute an approximation \( w_{t} \) of the
function \( w \) from below for which \( 1-\eps < \sum_{x} w_{t}(x) \) then we
have \( | w(x)-w_{t}(x) | < \eps \) for all \( x \). 
 \end{proof}

% Gacs source line 1743
 \begin{theorem}
  \lean{Kolmogorov.IsUniversalSemimeasure, Kolmogorov.exists_universalSemimeasure}
  \leanok
  \uses{def:semimeasure} \label{thm:univsem} 
There is a universal constructive semimeasure.
 \end{theorem}

% Gacs source line 1746
\begin{proof}
  \leanok
  \uses{def:semimeasure} 
 Every computable family \( w_{i} \) \( (i=1,2,\dotsc) \) of semimeasures is
dominated by the semimeasure \( \sum_{i} 2^{-i}w_{i} \). 
Since there are only
countably many constructive semimeasures there exists a semimeasure
dominating them all. 
The only point to prove is that the dominating
semimeasure can be made constructive. 

Below, we construct a function \( \mu_{p}(x) \) lower
semicomputable in \( p,x \) such 
that \( \mu_{p} \) as a function of \( x \) is a constructive semimeasure for
all \( p \) and all constructive semimeasures occur among the \( w_{p} \). 
Once we have \( \mu_{p} \) we are done because for any positive computable
function \( \dg(p) \) with \( \sum_{p} \dg(p) \le 1 \) the semimeasure 
\( \sum_{p} \dg(p) \mu_{p} \) is obviously lower
semicomputable and dominates all the
\( \mu_{p} \). 

Let \( S_{p}(x) \) be the lower semicomputable function with G\"odel number \( p \)
and let \( S_{p}^{t}(x) \) be a function recursive in \( p,t,x \) with rational
values, dondecreasing in \( t \), such that 
\( \lim_{t} S^{t}_{p}(x) = \max\{0,S_{p}(x)\} \) and for each \( t,p \), the function 
\( S_{p}^{t}(x) \) is
different from 0 only for \( x \le t \). Let us define \( \mu_{p}^{t} \)
recursively in \( t \) as follows. 
Let \( \mu_{p}^0(x) = 0 \). 
Suppose that \( \mu_{p}^{t} \) is already defined. 
If \( S_{p}^{t+1} \) is a semimeasure then we
define \( \mu_{p}^{t+1}=S_{p}^{t+1} \), otherwise \( \mu_{p}^{t+1} = \mu_{p}^{t} \).
Let \( \mu_{p} = \lim_{t} \mu_{p}^{t} \). 
The function \( \mu_{p}(x) \) is, by its
definition, lower semicomputable in \( p,x \) and a semimeasure for each fixed
\( p \). 
It is equal to \( S_{p} \) whenever \( S_{p} \) is a semimeasure, hence it
enumerates all constructive semimeasures.
 \end{proof}

% Gacs source line 1786
\begin{definition}
  \lean{Kolmogorov.IsUniversalSemimeasure}
  \leanok
  \uses{thm:univsem}\label{def:m}
We choose a fixed universal constructive semimeasure and call it \( \m(x) \).
\end{definition}

% Gacs source line 1793
 \begin{definition}
  \label{def:m-of-nu}
  \notready
For an arbitrary constructive semimeasure \( \nu \) let us define
 \begin{equation}\label{eq:m(nu)}
   \m(\nu) = 
\sum\setof{\m(p) : \txt{the (self-delimiting) program \( p \) computes \( \nu \)}}.
 \end{equation}
Similar notation will be used for other objects: for example
now \( \m(f) \) make sense for a recursive function \( f \).
 \end{definition}

% Gacs source line 1805
 \begin{theorem}
  \uses{def:m-of-nu, def:m}
  \notready\label{thm:exact-domin} 
For all constructive semimeasures \( \nu \) and for all strings \( x \), we have 
 \begin{equation}\label{eq:exact-domin}
  \m(\nu) \nu(x) \lem \m(x).
 \end{equation}
 \end{theorem}

\begin{formalization}
Not formalized. The library proves domination of every lower semicomputable semimeasure by the universal one, but not this sharpened form with the weight of the semimeasure itself as the constant.
\end{formalization}

% Gacs source line 1811
 \begin{proof}
  \uses{thm:univsem}
  \notready
Let us repeat the proof of Theorem~\ref{thm:univsem}, using \( \m(p) \) in place
of \( \dg(p) \).
We obtain a new constructive semimeasure \( \m'(x) \) with the property
\( \m(\nu)\nu(x) \le \m'(x) \) for all \( x \).
Noting \( \m'(x) \eqm \m(x) \) finishes the proof.
 \end{proof}

% Gacs source line 1826
 \begin{definition}
  \label{def:code}
  \uses{def:prefix-free}
A (binary) \df{code} is any
mapping from a set \( E \) of binary sequences to a set of objects. 
This mapping is the \df{decoding function} of the code. 
If the mapping is
one-to-one, then sometimes the set \( E \) itself is also called a code.

A code is a \df{prefix code} if \( E \) is a prefix-free set.
 \end{definition}

% Gacs source line 1838
\begin{definition}
  \label{def:first-shortest}
  \uses{def:K}
We call string \( p \) the first
shortest description of \( x \) if it is the lexicographically first
binary word \( q \) with \( |q| = C(x) \) and \( U(q)=x \). 
\end{definition}

% Gacs source line 1853
 \begin{lemma}[Kraft's Inequality, see~\protect\cite{CsiszarKornerBook81}]
  \lean{Kolmogorov.exists_prefixFree_code_of_kraft_le_one, Kolmogorov.finset_kraft_progWeight_le_one}
  \leanok
  \uses{def:prefix-free, def:code}
\label{l:kraft} 
For any sequence \( l_{1},l_{2},\ldots \) of
natural numbers, there is a prefix code with exactly these numbers as
codeword lengths, if and only if
 \begin{equation}	\sum_{i} 2^{-l_{i}} \le 1. \label{eq:kraft} 
 \end{equation}
 \end{lemma}

% Gacs source line 1861
\begin{proof}
  \leanok
  \uses{def:code} 
Recall the standard correspondence \( x \leftrightarrow [x] \) between
binary strings and binary subintervals of the interval \( \clint{0}{1} \).  
A prefix code corresponds to a set of disjoint intervals, and the 
length of the interval \( [x] \) is \( 2^{-\len{x}} \). 
This proves 
that~\eqref{eq:kraft} holds for the lengths of codewords of a prefix code.

Suppose that \( l_{i} \) is given and~\eqref{eq:kraft} holds. 
We can assume that the sequence \( l_{i} \) is nondecreasing. 
Chop disjoint, adjacent 
intervals \( I_{1},I_{2},\ldots \) of length \( 2^{-l_{1}},2^{-l_{2}},\ldots \) from
the left end of the interval \( [0,1] \). 
The right end of \( I_{k} \) is \( \sum_{j=1}^k 2^{-l_{j}} \). 
Since 
the sequence \( l_{j} \) is nondecreasing, 
all intervals \( I_{j} \) are binary intervals. 
Take the binary string  corresponding to \( I_{j} \) as the \( j \)-th codeword. 
 \end{proof}

% Gacs source line 1881
 \begin{corollary}
  \label{cor:m-le-K}
  \lean{Kolmogorov.complexityWeight_KP_le_aprioriMeasure, Kolmogorov.complexityWeight}
  \leanok
  \uses{l:kraft, def:K, def:m}
We have \( -\log \m(x) \lea K(x) \).
 \end{corollary}

% Gacs source line 1884
 \begin{proof}
  \proves{cor:m-le-K}
  \leanok
  \uses{l:kraft}
The lemma implies 
 \begin{equation} \sum_y 2^{-K(y \mvert x)} \le 1.  \label{eq:condHsum} 
 \end{equation}
  Since \( K(x) \) is an upper semicomputable function, the function
\( 2^{-K(x)} \) is lower semicomputable. 
Hence it is a constructive semimeasure, and we have \( -\log \m(x) \lea K(x) \).
 \end{proof}

% Gacs source line 1901
\begin{lemma}[Shannon-Fano code]
  \label{l:shannon-fano}
  \lean{Kolmogorov.exists_shannonFano_code}
  \leanok
  \uses{def:code, def:prefix-free} (see~\cite{CsiszarKornerBook81})
 Let \( w_{1},w_{2},\ldots \) be positive numbers with \( \sum_{j} w_{j} \le 1 \).
There is a binary prefix code \( p_{1},p_{2},\ldots \) where the codewords
are in lexicographical order, such that 
 \begin{equation} 	|p_{j}| \le -\log w_{j} + 2. \label{eq:shfano} 
 \end{equation}
 \end{lemma}

\begin{formalization}
The Lean version is finite and multiplicative: for finitely many weights it produces codeword lengths with $2^{-\len{p_j}} \le w_j < 2 \cdot 2^{-\len{p_j}}$. That is one bit sharper than the additive constant here, because we drop the requirement that the codewords follow the lexicographic order of the weights.
\end{formalization}

% Gacs source line 1908
 \begin{proof}
  \leanok
  \uses{def:code} We follow the construction above and cut off disjoint,
adjacent (not necessarily binary) intervals \( I_{j} \) of length \( w_{j} \) from 
the left end of \( [0,1] \). Let \( v_{j} \) be the length of the longest binary 
intervals contained in \( I_{j} \). 
Let \( p_{j} \) be the binary word corresponding to the first one of these. 
Four or fewer intervals of length \( v_{j} \) cover \( I_{j} \). 
Therefore~\eqref{eq:shfano} holds. 
\end{proof}

% Gacs source line 1924
\begin{theorem}[Coding Theorem, see~\protect\cite{Levin74,GacsSymm74,Chaitin75}]
  \lean{Kolmogorov.universalSemimeasure_equiv_prefixComplexity}
  \leanok
  \uses{def:m, def:K, l:kraft, l:shannon-fano}
\label{thm:coding} 
  \begin{equation}  	K(x) \eqa -\log\m(x) \label{eq:codingth} 
  \end{equation}
 \end{theorem}

% Gacs source line 1930
\begin{proof}
  \leanok
  \uses{l:shannon-fano, def:m} 
 We construct a self-delimiting p.r.~function \( F(p) \) with the
property that \( C_{F}(x) \le -\log\m(x) + 4 \). The function \( F \) to be 
constructed is the decoding function of a prefix code hence the 
code-construction of Lemma~\ref{l:kraft} proposes itself. 
But since
the function \( \m(x) \) is only lower semicomputable, it is given only by a
sequence converging to it from below. 

Let \( \setof{ \pair{z_{t}}{k_{t}} : t=1,2,\ldots } \) be a recursive enumeration
of the set \( \setof{ \pair{x}{k} : k < \m(x) } \) without repetition. Then 
 \[ 
  \sum_{t} 2^{-k_{t}} = \sum_{x} \sum_{z_{t}=1} 2^{-k_{t}} \le \sum_{x} 2\m(x) < 2.
 \]
  Let us cut off consecutive adjacent, disjoint intervals \( I_{t} \) of
length \( 2^{-k_{t}-1} \) from the left side of the interval \( \clint{0}{1} \). 
We define \( F \) as follows. If \( [p] \) is a largest binary subinterval of
some \( I_{t} \) then \( F(p)=z_{t} \). Otherwise \( F(p) \) is undefined. 

The function \( F \) is obviously self-delimiting and partial recursive. It 
follows from the construction that for every \( x \) there is a \( t \) with 
\( z_{t}=x \) and \( 0.5\m(x) < 2^{-k_{t}} \). Therefore, for every \( x \) 
there is a \( p \) such that \( F(p)=x \) and \( |p| \le -\log\m(x) + 4 \). 
 \end{proof}

% Gacs source line 1987
\begin{definition}
  \label{def:PT}
  \lean{Kolmogorov.aprioriMeasure}
  \leanok
  \uses{def:sd-machine}
Let \( P_{T}(x) \) be the probability that the self-delimiting
machine \( \cT \) gives out result \( x \) after receiving random bits as inputs.
\end{definition}

% Gacs source line 2012
 \begin{definition}
  \label{def:m-is-PT}
  \lean{Kolmogorov.aprioriMeasure}
  \leanok
  \uses{def:PT, def:m}From now on let us define
 \begin{equation}\label{eq:m(x)}
   \m(x) = P_{T}(x).
 \end{equation}
 \end{definition}

%%%% ===================== SECTION 1.7 The statistics of description length =====================


\section{The statistics of description length }
How many descriptions does an object have? The coding theorem answers several questions
of this kind at once. This section is where the formalization does the most independent
work, and where our statements depart most often from the printed ones.


% Gacs source line 2026
\begin{theorem}
  \lean{Kolmogorov.progCount, Kolmogorov.progCount_le_aprioriMeasure_pair, Kolmogorov.inv_two_pow_mul_progCount_le_complexityWeight_pair, Kolmogorov.aprioriMeasure_pair_le_progCount}
  \leanok
  \uses{thm:sdopt, def:m, def:pair-complexity} \label{thm:Hstat} 
  Let $f(x,n)$ be the number of binary strings $p$ of length $n$ with
$T(p)=x$ for the universal s.d.~interpreter $T$ defined in 
Theorem~\ref{thm:sdopt}.
  Then for every $n \ge K(x)$, we have
  \begin{equation}
	\log f(x,n) \eqa n - K(x,n). \label{eq:Hstat} 
  \end{equation}
 \end{theorem}

\begin{formalization}
Formalized in the multiplicative form $2^{-n} f(x,n) \eqm \m(x,n)$ that the proof below actually establishes, in $[0,\infty]$ rather than through logarithms. The two halves are separate results, and the lower half carries the hypothesis $n \ge K(x,n) + c$, which is what the padding argument consumes. That hypothesis is stronger than the printed $n \ge K(x)$ by up to a logarithmic term; in the gap between them the printed claim needs $x$ to have a program of length exactly $n$, which holds at $n = K(x)$ but is not proved in the interior.
\end{formalization}

% Gacs source line 2040
  \begin{corollary}
  \label{cor:shortest-descriptions}
  \lean{Kolmogorov.progCount_shortest_le}
  \leanok
  \uses{thm:Hstat} The number of shortest descriptions of any object is
bounded by a universal constant.
 \end{corollary}

% Gacs source line 2052
\begin{proof}[Proof of Theorem \protect\ref{thm:Hstat}]
  \proves{thm:Hstat}
  \leanok
  \uses{thm:coding, def:m}
 It is more convenient to prove equation~\eqref{eq:Hstat}
in the form 
 \begin{equation} \label{eq:Hstatmult} 
 	2^{-n}f(x,n) \eqm \m(x,n)
 \end{equation}
 where $\m(x,y) = \m(\ang{ x,y} )$.
First we prove $\lem$.
Let us define a p.r.~self-delimiting function $F$ by
 $F(p) = \tupcod{T(p), \len{p}}$.
Applying $F$ to a coin-tossing argument, $2^{-n}f(x,n)$ is the
probability that the pair $\ang{x, n} $ is obtained.
Therefore the left-hand side of~\eqref{eq:Hstatmult}, as a function of
$\ang{ x,n}$, is a constructive semimeasure, dominated by $\m(x,n)$.

  Now we prove $\gem$.
  We define a self-delimiting p.r.~function $G$ as follows.
  The machine computing $G(p)$ tries to decompose $p$ into
three segments $p=\bg(c)^o v w$ in such a way that $T(v)$ is a pair 
$\ang{ x, \len{p} + c }$.
  If it succeeds then it outputs $x$.
  By the universality of $T$, there is a binary string $g$ such that
$T(gp)=G(p)$ for all $p$.
  Let $r=\len{g}$.
 For an arbitrary pair $x,n$, let $q$ be a shortest binary string with
$T(q) = \tupcod{ x,n}$, and $w$ an arbitrary string of length
 \[
 	l = n - \len{q} - r - \len{\bg(r)^o}.
 \]
  Then $G(\bg(r)^o q w) = T(g \bg(r)^o q w) = x$.
  Since $w$ is arbitrary here, there are $2^l$ binary strings $p$ of length
$n$ with $T(p) = x$.
 \end{proof}

% Gacs source line 2089
 \begin{definition}
  \label{def:Dn}
  \uses{def:K}
  Let $g_{T}(n)$ be the number of objects $x \in \dS$ with $K(x)=n$, and $D_{n}$
the set of binary strings $p$ of length $n$ for which $T(p)$ is
defined.
  Let us define the moving average
 \[
	h_{T}(n,c) = \frac{1}{2c+1}\sum_{i = -c}^c g_{T}(n+i).
 \]
 \end{definition}

% Gacs source line 2101
\begin{theorem}
  \lean{Kolmogorov.haltCount_le_complexityWeight, Kolmogorov.complexityWeight_le_haltCount, Kolmogorov.windowCount_eq_complexityWeight}
  \leanok
  \uses{def:Dn, thm:Hstat} \label{thm:exactstat} (\cite{SolovayManu}) 
 There is a natural number $c$ such that 
 \begin{align}
	\log\card{D_{n}}	&\eqa n-K(n),	\label{eq:Dn}
  \\	\log h_{T}(n,c)	&\eqa n-K(n).	\label{eq:aver}
 \end{align}
 \end{theorem}

\begin{formalization}
Both halves are formalized multiplicatively. The bound on $\card{D_n}$ from below carries the guard $K(n) + c \le n$, and the moving-average half is stated as an equality of window counts with the complexity weight.
\end{formalization}

% Gacs source line 2124
\begin{lemma}
  \lean{Kolmogorov.pairMarginal_eq_aprioriMeasure}
  \leanok
  \uses{def:m} \label{l:projmeas} 
 \[
 	\sum_y \m(x,y) \eqm \m(x).
 \]
 \end{lemma}

% Gacs source line 2130
\begin{proof}
  \leanok
  \uses{def:m, thm:univsem}
  The left-hand side is a constructive semimeasure therefore it is
dominated by the right side.
  To show $\gem$, note that by equation~\eqref{eq:xx.x} as applied to $H$, we
have $K(x,x) \eqa K(x)$.
  Therefore $\m(x) \eqm \m(x,x) < \sum_y \m(x,y)$.
 \end{proof}

% Gacs source line 2138
\begin{proof}[Proof of Theorem \protect\ref{thm:exactstat}:]
  \proves{thm:exactstat}
  \leanok
  \uses{l:projmeas, thm:Hstat}
 Let $d_{n} = \card{D_{n}}$. Using Lemma~\ref{l:projmeas} and 
Theorem~\ref{thm:Hstat} we have 
 \[
	g_{T}(n) \le d_{n} = \sum_{x} f(x,n) \eqm 2^{n} \sum_{x} \m(x,n) 
			\eqm 2^{n} \m(n).
 \]
  Since the complexity $H$ has the same continuity 
property~\eqref{eq:continui} as $K$, we can derive from here
 $h_{T}(n,c) \le 2^{n}\m(n)O(2^c c^2)$.
  To prove $h_{T}(n,c) \gem d_{n}$ for an appropriate $c$, we will prove that
the complexity of at least half of the elements of $D_{n}$ is near $n$.

  For some constant $c_{0}$, we have $K(p) \le n + c_{0}$ for all $p \in
D_{n}$. 
  Indeed, if the s.d.~p.r.~interpreter $F(p)$ is equal to $p$
where $T(p)$ is defined and is undefined otherwise then $C_{F}(p) \le
n$ for $p \in D_{n}$.
  By $K(p,\len{p})\eqa K(p)$, and a variant of Lemma~\ref{l:projmeas}, we have
 \[
 \sum_{p \in D_{n}}\m(p) \eqm \sum_{p \in D_{n}} \m(p,n) \eqm \m(n).
 \]
  Using the expression~\eqref{eq:Dn} for $d_{n}$, we see that there is a
constant $c_{1}$ such that 
 \[
  d_{n}^{-1} \sum_{p \in D_{n}} 2^{-K(p)} \le 2^{-n+c_{1}}.
 \]
  Using Markov's Inequality, hence the number of elements $p$ of $D_{n}$
with $K(p) \ge n-c_{1}-1$ is at least $d_{n}/2$.
  We finish the proof making $c$ to be the maximium of $c_{0}$ and $c_{1}+1$.
 \end{proof}

% Gacs source line 2189
 \begin{definition}
  \label{def:Kplus}
  \lean{Kolmogorov.KPlus}
  \leanok
  \uses{def:K}
Let us introduce the following function of natural numbers:
 \[	
  K^{+}(n) = \max_{k \le n} K(k).
 \]
 \end{definition}

% Gacs source line 2202
\begin{theorem}
  \lean{Kolmogorov.KPlus_eq_size_add}
  \leanok
  \uses{def:Kplus} \label{thm:monbd} 
We have $K^+(n) \eqa \log n + K(\flo{ \log n })$.
  \end{theorem}

% Gacs source line 2205
\begin{proof}
  \proves{thm:monbd}
  \leanok
  \uses{def:Kplus, def:sd-machine} 
 To prove $\lea$, we construct a s.d.~interpreter as follows.
  The machine computing $F(p)$ finds a decomposition $uv$ of $p$ such that
$T(u)=\len{v}$, then outputs the number whose binary representation (with
possible leading zeros) is the binary string $v$.
  With this $F$, we have
 \begin{equation} \label{eq:monbd1}   
   K(k) \lea C_{F}(k) \le \log n + H (\flo{ \log n }) 
 \end{equation}
  for all $k \le n$.
 The numbers between $n/2$ and $n$ are the ones whose binary representation
has length $\flo{ \log n} + 1$.
  Therefore if the bound~\eqref{eq:monbd} is sharp for most $x$ then the 
bound~\eqref{eq:monbd1} is sharp for most $k$.
 \end{proof}

% Gacs source line 2231
\begin{theorem}[see~\protect\cite{SolovayManu}]
  \lean{Kolmogorov.exists_computable_upper_bound_eq_infinitely_often}
  \leanok
  \uses{def:K, thm:monbd} \label{thm:infsharp} 
  There is a computable upper bound $G(n)$ of the function $K(n)$ with the
property that $G(n)=K(n)$ holds for infinitely many $n$.
 \end{theorem}

% Gacs source line 2241
 \begin{definition}
  \label{def:K-with-time}
  \uses{def:K}
We know that for some constant $c$, the
function $K(n)$ is bounded by $2\log n + c$.
Let us fix such a $c$.

For any number $n$ and s.d.~ Turing machine $\cM$, let $K_{\cM}(n;t)$ be
the minimum of $2\log n + c$ and the lengths of descriptions from which
$\cM$ computes $n$ in $t$ or fewer steps.
 \end{definition}

% Gacs source line 2255
 \begin{definition}
  \label{def:universal-machine-V}
  \uses{def:K-with-time}
Let us agree now that the universal p.r.~function
is computed by a universal Turing
machine $\cV$ with the property that for any Turing machine $\cM$ there is
a constant $m$ such that the number of steps needed to simulate $t$ steps
of $\cM$ takes no more than $mt$ steps on $\cV$.
Let $\cT$ be the Turing machine constructed from $\cV$ computing the
optimal s.d.~interpreter $T$, and let us write 
 \begin{align*}
  K(n,t)=K_{\cT}(n,t).
 \end{align*}
 \end{definition}

% Gacs source line 2275
\begin{proof}[Proof of Theorem \protect\ref{thm:infsharp}]
  \proves{thm:infsharp}
  \leanok
  \uses{def:K-with-time, def:universal-machine-V}
 First we prove that there are constants $c_{0},c_{1}$ such that $K(n;t) <
K(n;t-1)$ implies 
 \begin{equation} \label{eq:timetoo} 
 	K(\ang{ n,t} ; c_{0} t) \le K(n;t) + c_{1}.
 \end{equation}
  Indeed, we can construct a Turing machine $\cM$ simulating the 
work of $\cT$ such that if $\cT$ outputs a number $n$ in $t$
steps then $\cM$ outputs the pair $\ang{n,t}$ in $2t$ steps:
 \[
	K_{\cM}(\tupcod{ n,t} ; 2t) \le K(n,t).
 \]
  Combining this with the inequality~\eqref{eq:fastopt} gives the
inequality~\eqref{eq:timetoo}.
  We define now a recursive function $F$ as follows.
  To compute $F(x)$, the s.d.~Turing machine first tries
to find $n,t$ such that $x = \ang{ n,t}$.
 (It stops even if it did not find them.)
  Then it outputs $K(x; c_{0} t)$.
  Since $F(x)$ is the length of some description of $x$, it is an upper
bound for $K(x)$.
  On the other hand, suppose that $x$ is one of the infinitely many
integers of the form $\ang{ n,t } $ with
 \[
	K(n) = K(n;t) < K(n; t-1). 
 \]
  Then the inequality~\eqref{eq:timetoo} implies $F(x) \le K(n) + c_{1}$ while
$K(n) \lea K(x)$ is known (it is the equivalent of~\ref{thm:xKx} for $H$), so
$F(x) \le K(n) + c$ for some constant $c$. Now if $F$ would not be equal to
$H$ at infinitely many places, only close to it (closer than $c$) then we
can decrease it by some additive constant less than $c$ and modify it at
finitely many places to obtain the desired function $G$.
 \end{proof}


\bibliographystyle{alpha}
\bibliography{blueprint}
